\documentclass[../../main.tex]{subfiles}
\graphicspath{{\subfix{../../image/}}}

\begin{document}

\section{习题}\label{section:ch09-7-习题}

在下列习题中，$U$ 总表示 $\mathbb{R}^n$ 中边界光滑的有界开集。

\begin{exercise}\label{exercise:ch09-7-1}
设向量场 $\mathbf{v}$ 光滑。对这个情形，通过求解 ODE
\begin{align*}
\begin{cases}\dot{\mathbf{x}}(t)=-\mathbf{v}(\mathbf{x}(t))&(t\geqslant 0),\\ \mathbf{x}(0)=y\end{cases}
\end{align*}
来给出 \S9.1 引理的另一证明。把解写为 $\mathbf{x}(t,y)$ 以显示对初值 $y$ 的依赖。对每个固定时刻 $t>0$，映射 $y\mapsto\mathbf{x}(t,y)$ 连续，故有不动点。推得 $\mathbf{v}$ 在闭球 $B(0,r)$ 内有零点。
\end{exercise}

\begin{exercise}\label{exercise:ch09-7-2}
设 $a:\mathbb{R}\to\mathbb{R}$ 连续，且每当 $f_n\rightharpoonup f$ 弱于 $L^2(0,1)$ 时，$a(f_n)\rightharpoonup a(f)$ 弱于 $L^2(0,1)$。证明 $a$ 是仿射函数；即 $a$ 具有形式
\begin{align*}
a(z)=\alpha z+\beta\qquad (z\in\mathbb{R})
\end{align*}
对常数 $\alpha,\beta$。
\end{exercise}

\begin{exercise}[罚方法]\label{exercise:ch09-7-3}设 $\varepsilon>0$。定义
\begin{align*}
\beta_\varepsilon(z)\triangleq\begin{cases}0&\text{若 }z\geqslant 0,\\ \dfrac z\varepsilon&\text{若 }z\leqslant 0,\end{cases}
\end{align*}
并设 $u_\varepsilon\in H_0^1(U)$ 是
\begin{align}
\begin{cases}
-\Delta u_\varepsilon+\beta_\varepsilon(u_\varepsilon)=f&\text{在 }U\text{ 中},\\
u_\varepsilon=0&\text{在 }\partial U\text{ 上},
\end{cases}\label{exercise:ch09-7-3-star}
\end{align}
的弱解，其中 $f\in L^2(U)$。证明当 $\varepsilon\to 0$ 时，$u_\varepsilon\rightharpoonup u$ 弱于 $H_0^1(U)$，其中 $u$ 是变分不等式
\begin{align*}
\int_U Du\cdot D(w-u)\,\mathrm{d}x\geqslant\int_U f(w-u)\,\mathrm{d}x
\end{align*}
对所有 $w\in H_0^1(U)$、$w\geqslant 0$ a.e. 的唯一非负解。用\eqref{exercise:ch09-7-3-star} 逼近变分不等式即罚方法。
\end{exercise}

\begin{exercise}[时间周期解]\label{exercise:ch09-7-4}设
\begin{align*}
\begin{cases}
u_t-\Delta u=f&\text{在 }U\times(0,\infty)\text{ 中},\\
u=0&\text{在 }\partial U\times(0,\infty)\text{ 上},\\
u=g&\text{在 }U\times\{t=0\}\text{ 上},
\end{cases}
\end{align*}
其中 $g\in L^2(U)$，$f\in L^\infty(U\times(0,\infty))$。设 $\tau>0$，$f$ 关于 $t$ 是 $\tau$-周期的；即 $f(x,t)=f(x,t+\tau)$（$x\in U$，$t\geqslant 0$）。证明存在唯一函数 $g\in L^2(U)$，使相应解 $u$ 是 $\tau$-周期的。
\end{exercise}

\begin{exercise}\label{exercise:ch09-7-5}
考察非线性边值问题
\begin{align*}
\begin{cases}
-\Delta u+b(Du)=f&\text{在 }U\text{ 中},\\
u=0&\text{在 }\partial U\text{ 上}.
\end{cases}
\end{align*}
用 Banach 不动点定理证明：若 $b:\mathbb{R}^n\to\mathbb{R}$ Lipschitz 连续且 $\operatorname{Lip}(b)$ 充分小，则存在唯一弱解 $u\in H^2(U)\cap H_0^1(U)$。
\end{exercise}

\begin{exercise}\label{exercise:ch09-7-6}
设 $f:\mathbb{R}\to\mathbb{R}$ Lipschitz 连续、有界，$f(0)=0$，$f'(0)>\lambda_1$，$\lambda_1$ 表示 $-\Delta$ 在 $H_0^1(U)$ 上的主特征值。用上下解方法证明存在弱解 $u$ 满足
\begin{align*}
\begin{cases}
-\Delta u=f(u)&\text{在 }U\text{ 中},\\
u=0&\text{在 }\partial U\text{ 上},\\
u>0&\text{在 }U\text{ 中}.
\end{cases}
\end{align*}
\end{exercise}

\begin{exercise}\label{exercise:ch09-7-7}
设 $\underline{u}$、$\overline{u}$ 是 \S9.3 边值问题\eqref{align:ch09-3-eq01} 的光滑下解与上解。用极值原理直接验证
\begin{align*}
\underline{u}=u_0\leqslant u_1\leqslant\cdots\leqslant u_k\leqslant\cdots\leqslant\overline{u},
\end{align*}
其中 $\{u_k\}_{k=0}^\infty$ 如 \S9.3 所定义。
\end{exercise}

\begin{exercise}[解的非紧族]\label{exercise:ch09-7-8}
(a) 设 $n\geqslant 3$。求常数 $c$ 使
\begin{align*}
u(x)\triangleq(1+|x|^2)^{\frac{2-n}{2}}
\end{align*}
求解 Yamabe 方程
\begin{align*}
-\Delta u=c\,u^{\frac{n+2}{n-2}}\quad\text{在 }\mathbb{R}^n\text{ 中}.
\end{align*}
注意临界指数 $\frac{n+2}{n-2}$ 的出现。（与\cref{exercise:ch04-06} 比较。）
(b) 验证对每个 $\lambda>0$，
\begin{align*}
u_\lambda(x)\triangleq\left(\frac\lambda{\lambda^2+|x|^2}\right)^{\frac{n-2}{2}}
\end{align*}
也是解。
(c) 证明对每个 $\lambda$，
\begin{align*}
\|u_\lambda\|_{L^{2^*}(\mathbb{R}^n)}=\|u\|_{L^{2^*}(\mathbb{R}^n)},\qquad \|Du_\lambda\|_{L^2(\mathbb{R}^n)}=\|Du\|_{L^2(\mathbb{R}^n)},
\end{align*}
从而 $\{u_\lambda\}_{\lambda>0}$ 在 $L^{2^*}(\mathbb{R}^n)$ 中不预紧。
\end{exercise}

\begin{exercise}\label{exercise:ch09-7-9}
设 $u$ 求解半线性热方程
\begin{align*}
u_t-\Delta u=f(u)\quad\text{在 }\mathbb{R}^n\times(0,\infty)\text{ 中}.
\end{align*}
假设 $u$ 及其导数当 $|x|\to\infty$ 时快速趋于零。
(a) 证明
\begin{align*}
\frac{\mathrm{d}}{\mathrm{d}t}\int_{\mathbb{R}^n}\frac12|Du|^2-F(u)\,\mathrm{d}x=-\int_{\mathbb{R}^n}u_t^2\,\mathrm{d}x,
\end{align*}
其中 $F'=f$，$F(0)=0$。
(b) 现证明
\begin{align*}
\frac{\mathrm{d}}{\mathrm{d}t}\int_{\mathbb{R}^n}|x|^2\left(\frac12|Du|^2-F(u)\right)\,\mathrm{d}x=-\int_{\mathbb{R}^n}u_t^2|x|^2-2nF(u)+(n-2)|Du|^2\,\mathrm{d}x.
\end{align*}
这是 Derrick--Pohozaev 恒等式的抛物类比。（Bauman--Chen--Phillips, Euro.\ J.\ Applied Math.\ 6 (1995), 115--126）
\end{exercise}

\begin{exercise}\label{exercise:ch09-7-10}
设 $K\subset\mathbb{R}^n$ 为闭、凸、非空集合。定义
\begin{align*}
I[x]\triangleq\begin{cases}0&\text{若 }x\in K,\\ \infty&\text{若 }x\notin K.\end{cases}
\end{align*}
用 $K$ 的几何显式确定 $A=\partial I$、$J_\lambda=(I+\lambda A)^{-1}$、$A_\lambda=\frac{I-J_\lambda}{\lambda}$（$\lambda>0$）。
\end{exercise}

\begin{exercise}\label{exercise:ch09-7-11}
给出一个简单例子表明流
\begin{align}
\mathbf{u}'\in-\partial I[\mathbf{u}]\quad (t>0)\label{exercise:ch09-7-11-star}
\end{align}
可能不可逆。（即找一个 Hilbert 空间 $H$ 及凸、真、下半连续函数 $I:H\to(-\infty,+\infty]$，使\eqref{exercise:ch09-7-11-star} 的半群解满足 $S(t)u=S(t)\widehat{u}$，对某 $t>0$ 及 $u\neq\widehat{u}$。）
\end{exercise}

\begin{exercise}\label{exercise:ch09-7-12}
设 $u=u(x,t)$ 表示沙堆在 $x\in\mathbb{R}^2$ 处的高度，沙堆以速率 $f=f(x,t)\geqslant 0$ 加沙增长。假设稳定条件
\begin{align*}
|Du|\leqslant 1,
\end{align*}
即沙堆任何地方斜率不超过 1。如常 $D_xu=Du$。提议动力学
\begin{align*}
u_t-\operatorname{div}(aDu)=f\quad\text{在 }\mathbb{R}^2\times(0,\infty)\text{ 中},
\end{align*}
其中 $a=a(x,t)\geqslant 0$ 描述沙向下滚落的流速，即沿 $-Du$ 方向。最后假设
\begin{align*}
\operatorname{spt}a\subseteq\{|Du|=1\},
\end{align*}
即沙只在斜率为一时才向下流。证明上述蕴涵
\begin{align*}
f-u_t\in\partial I[u],
\end{align*}
其中凸函数
\begin{align*}
I[u]\triangleq\begin{cases}0&\text{若 }u\in L^2(\mathbb{R}^2),\ |Du|\leqslant 1\text{ a.e.},\\ \infty&\text{否则}.\end{cases}
\end{align*}
\end{exercise}

\begin{exercise}\label{exercise:ch09-7-13}
设 $u$ 是 \S9.6.3 梯度流方程组\eqref{align:ch09-6-eq33} 的光滑解，其中 $L$ 满足一致凸性条件\eqref{align:ch09-6-eq29}。证明存在常数 $C,\gamma>0$ 使
\begin{align*}
\int_U u_t^2(x,t)\,\mathrm{d}x\leqslant Ce^{-\gamma t}\qquad (t>0).
\end{align*}
\end{exercise}

\begin{exercise}\label{exercise:ch09-7-14}
设 $u$ 是非线性热方程
\begin{align*}
u_t-\Delta\phi(u)=0\quad\text{在 }\mathbb{R}^n\times(0,\infty)\text{ 中}
\end{align*}
的光滑解，其中 $\phi'>0$。假设 $|Du|>0$，推导公式
\begin{align*}
\frac{\mathrm{d}}{\mathrm{d}t}\int_{\mathbb{R}^n}|Du|\,\mathrm{d}x=-\int_{\mathbb{R}^n}\frac{|D^2v|^2}{|Dv|}-\frac{|D^2vDv|^2}{|Dv|^3}\,\mathrm{d}x\leqslant 0,
\end{align*}
其中 $v\triangleq\phi(u)$。
\end{exercise}

\end{document}
